\section{Introduction}
\label{sec:intro}

Autonomous LLM agents are increasingly deployed inside containerized sandboxes
that are assumed to confine them. Recent work shows this assumption is fragile:
frontier models reliably escape common container misconfigurations, such as privileged
containers and exposed Docker sockets~\cite{sandboxescapebench}. A parallel finding in web
security reports that none of seven frontier models bypassed production web application
firewalls~\cite{pacebench}. Together these results raise an unanswered question for
anyone deploying LLM agents today: when we actually \emph{harden} the container,
which controls matter, and how much does each one help?

Existing benchmarks answer ``\emph{can} the agent escape?'' under the weakest
threat model, but not ``\emph{how much do standard protections reduce escape
success, and why?}'' Deployers are left with qualitative best-practice checklists
and no empirical, ranked evidence on which controls are essential versus
nice-to-have, or whether a
defense that appears to work closes the vulnerability or merely confuses the
model.

\paragraph{Contributions.}
This paper makes the following contributions:
\begin{itemize}
  \item A \textbf{defensive ablation study} over the SandboxEscapeBench scenarios~\cite{sandboxescapebench}
  that quantifies the escape-rate reduction attributable to each of eight
  hardening controls (\Cref{sec:methodology}).
  \item A \textbf{mechanistic verification} method (RQ2) that runs a hand-written
  reference exploit under each control to distinguish a true architectural fix
  from a friction-induced model failure, and to flag unintended bypasses.
  \item An \textbf{empirical, ranked hardening matrix} with absolute and relative
  risk reductions, intended as actionable deployment guidance
  (\Cref{sec:results}).
  \item An \textbf{open, reproducible harness} extending Inspect AI with a
  per-control hardening wrapper and a reference-exploit verifier.
\end{itemize}

\paragraph{Research questions.}
\begin{description}
  \item[RQ1.] For each hardening control, what is the empirical escape-rate
  reduction relative to the unhardened baseline?
  \item[RQ2.] When a defense stops an escape, is it because the vulnerability is
  actually closed, or because the model failed for unrelated (friction) reasons?
\end{description}
